
\chapter{Diagrama de Objetos}\label{ch:diagramaobjetos}

El diagrama de objetos complementa al diagrama de clases conceptuales mostrando una \emph{instantánea} (snapshot) del sistema en un momento determinado: instancias concretas de las clases con valores reales en sus atributos y enlaces entre ellas. Mientras el diagrama de clases describe qué \emph{puede} existir, el diagrama de objetos muestra qué \emph{existe} en un escenario específico, lo que ayuda a validar el modelo conceptual.

\section{Escenario representado}

Se modela una instantánea realista del sistema centrada en una paciente, \textit{Ana López}, en tratamiento por un esguince de ligamento colateral medial de rodilla derecha. La instantánea incluye:

\begin{itemize}
  \item La paciente (\codeinline{ana}) atendida por el terapeuta \codeinline{drMartinez}.
  \item Su \codeinline{HistorialClinico} con un diagnóstico CIE-10 \codeinline{S83.4} y una \codeinline{NotaClinica} de evolución.
  \item Un \codeinline{PlanTratamiento} activo (\textit{Recuperación LCM rodilla derecha}) con la \codeinline{FaseTratamiento} \textit{Fase 2: Fortalecimiento} en curso.
  \item Una \codeinline{Cita} de revaloración programada en el \codeinline{Consultorio 3}, con una \codeinline{Notificacion} de recordatorio enviada a la paciente.
  \item Un \codeinline{DocumentoClinico} (radiografía) adjunto al expediente.
  \item La \codeinline{AsignacionRutina} de la rutina \textit{Fortalecimiento rodilla nivel 2}, con una \codeinline{SesionEjecucion} completada y un \codeinline{Feedback} cuya transcripción y resumen de IA se almacenan en línea, junto con un \codeinline{PuntoDolor} en la rodilla y una \codeinline{EscalaResultado} de rango de movimiento.
\end{itemize}

\begin{figure}[H]
  \centering
  \includegraphics[width=\textwidth,keepaspectratio]{FisioManager_DiagramaObjetos.png}
  \caption{Diagrama de Objetos de FisioManager: instantánea del tratamiento de rodilla de la paciente Ana López.}
  \label{fig:diagrama-objetos}
\end{figure}

\section{Propósito del diagrama}

Esta instantánea permite verificar varios aspectos del modelo conceptual:

\begin{itemize}
  \item \textbf{Coherencia de cardinalidades:} una paciente tiene exactamente un historial clínico y un terapeuta asignado; el plan de tratamiento agrupa fases y asignaciones de rutina.
  \item \textbf{Trazabilidad clínica:} cada nota clínica está vinculada al historial; cada feedback está vinculado a su sesión e incorpora su transcripción y resumen de IA, además de los puntos de dolor reportados.
  \item \textbf{Validación del paquete de agenda:} la cita se relaciona simultáneamente con paciente, terapeuta, sala y plan, y genera una notificación, mostrando que el paquete cubre todos los puntos de integración previstos.
  \item \textbf{Estado del sistema en un instante:} a diferencia del diagrama de clases, aquí cada objeto tiene valores concretos, lo que sirve como caso de prueba conceptual y guía para los datos de \textit{seed} en pruebas e historias de usuario.
\end{itemize}

\section{Código PlantUML del diagrama}

El código fuente del diagrama de objetos se entrega como \codeinline{07\_diagrama\_objetos.puml} en el directorio \codeinline{puml/}. Se regenera con:

\begin{lstlisting}[language=bash]
plantuml -tpng 07_diagrama_objetos.puml
\end{lstlisting}

El código fuente completo se incluye en el anexo de código fuente PlantUML del Modelo de Análisis.
